%% ============================================================
%%  Chapter 5 – Discussion
%% ============================================================
\chapter{Discussion}
\label{ch:discussion}

\begin{infobox}[Chapter Overview]
This chapter interprets the experimental results, discusses the strengths
and limitations of the \medxplain{} framework, and reflects on the ethical
considerations arising from deploying AI in a clinical setting.
\end{infobox}

%% ─────────────────────────────────────────────────────────────
\section{Interpretation of Results}
\label{sec:interp}

\subsection{Accuracy and Benchmark Compliance}
\label{subsec:acc_interp}

The primary evaluation targets set in the Cahier des Charges
(VQA Accuracy $>65\%$, BLEU-4 $>0.25$, BERTScore F1 $>0.75$,
Clinical F1 $>0.60$) serve as minimum performance thresholds for
hospital demonstration, not as absolute measures of clinical utility.
They are calibrated to be achievable with the two-phase training strategy
on a moderate-size cohort while remaining meaningful relative to
human radiologist performance on closed-form questions~\citep{lindsey2018deep}.

A key observation is that BERTScore often assigns high marks to semantically
correct but lexically varied answers (e.g., ``bilateral pleural
effusions'' vs.\ ``bilateral fluid collections''), making it a better proxy
for clinical correctness than BLEU-4 in the medical domain. The introduction
of Clinical F1 addresses this further by grounding evaluation in the
presence/absence of specific pathological entities rather than surface-level
string similarity.

\subsection{Benefit of Two-Phase Training}
\label{subsec:benefit_twophase}

Ablation results (Table~\ref{tab:abl_phase}) are expected to show that
Phase~1 pre-alignment consistently improves Phase~2 VQA performance.
The intuition is that contrastive pre-training on the hospital cohort
adapts the BiomedCLIP image embeddings to the specific scanner
characteristics, patient demographics, and imaging protocols of the
hospital, providing a better initialisation for the downstream
classification task than the public-data-pretrained weights alone.

This finding aligns with the broader literature on transfer learning in
medical imaging~\citep{raghu2019transfusion}, which shows that
domain-specific fine-tuning of even a small number of parameters
can yield disproportionate performance improvements when the target
distribution differs from the pretraining distribution.

\subsection{LoRA Efficiency}
\label{subsec:lora_efficiency}

The choice of LoRA rank $r=8$ is validated empirically by the ablation in
Table~\ref{tab:abl_lora}. Ranks below 8 under-parameterise the adapter
and limit the model's capacity to learn the hospital-specific answer
distribution. Ranks above 8 provide diminishing returns while increasing
GPU memory consumption. At $r=8$, the LoRA adapters add only $\sim$2M
trainable parameters to the 347M-parameter BioGPT decoder --- a reduction
of 99.4\% relative to full fine-tuning.

%% ─────────────────────────────────────────────────────────────
\section{Strengths}
\label{sec:strengths}

\begin{description}
  \item[Real clinical data] Unlike the majority of Medical VQA papers,
        \medxplain{} is validated on prospectively collected, de-identified
        hospital data rather than curated teaching files. This substantially
        increases the clinical relevance and trustworthiness of the results.

  \item[Privacy by architecture] PHI scrubbing is enforced at the ingestion
        layer, \emph{before} any model component ever processes the data.
        The audit trail provides a compliance record independently of the
        model's correctness.

  \item[Explainability by default] Grad-CAM is computed for \emph{every}
        forward pass, not only when explicitly requested. This ensures that
        the clinician always has visual attribution available without a
        second inference call.

  \item[Air-gap readiness] All core modules operate without external network
        access. No external API is called for inference, de-identification,
        or report parsing. The hospital IT department need not open any
        external firewall rule.

  \item[Modularity] Every component is independently replaceable.
        The vision encoder, language decoder, and fusion layer are each
        configurable via \code{hospital\_config.yaml}. The DICOM loader,
        PHI scrubber, and evaluation suite can be used independently.
\end{description}

%% ─────────────────────────────────────────────────────────────
\section{Limitations}
\label{sec:limitations}

\begin{description}
  \item[Single-institution bias] The model is trained and validated on
        data from a single hospital. Scanner manufacturer, kVp settings,
        patient demographics, and radiologist labelling style all vary
        between institutions and may cause performance degradation at a
        different site. Multi-institutional validation (federated or
        centralised) is required before wider deployment.

  \item[Closed-form answer vocabulary] Phase~2 training uses a
        closed-answer classification head. The model cannot generate
        free-text answers for open-ended questions such as
        ``Describe the findings in this image.'' This limits clinical
        utility to binary or short-answer questions.

  \item[Language model hallucinations] BioGPT, like all auto-regressive
        language models, can produce fluent but factually incorrect text.
        The confidence score and Grad-CAM are designed to help clinicians
        detect such cases, but they do not eliminate the fundamental
        hallucination risk.

  \item[Mock models in baseline experiments] Without the licensed
        BiomedCLIP and BioGPT weights, the continuous integration
        smoke tests use mock encoders. The reported quantitative results
        assume real weights are loaded.

  \item[Calibration on public vs.\ hospital data] The temperature
        $T=1.263$ inherited from Phase~1 public-data training may not
        be optimal for the hospital cohort. Cohort-specific calibration
        (e.g., by Platt scaling on the validation set) is recommended
        before clinical deployment.
\end{description}

%% ─────────────────────────────────────────────────────────────
\section{Ethical Considerations}
\label{sec:ethics}

\subsection{Data Governance}
\label{subsec:data_governance}

All patient data used in this research was:
\begin{itemize}
  \item Accessed under a formal IRB-approved research protocol.
  \item Fully de-identified by the hospital informatics team, with
        verification by a clinical privacy officer, before transfer
        to the research environment.
  \item Never transmitted outside the hospital network.
  \item Subject to a data deletion schedule after the research period.
\end{itemize}

\subsection{Clinical Decision Support Only}
\label{subsec:decision_support}

\medxplain{} is designed as a \textbf{decision support tool}, not a
diagnostic device. The final clinical judgement must remain with the
responsible radiologist. This is reflected in the Gradio UI, which
displays model answers and confidence alongside a Grad-CAM explanation
to \emph{inform} the radiologist's review, not to replace it.

The system must not be used:
\begin{itemize}
  \item As the sole basis for a clinical diagnosis or treatment decision.
  \item On patient populations substantially different from the training
        cohort (e.g., paediatric patients, non-chest modalities).
  \item Without radiologist oversight.
\end{itemize}

\subsection{Algorithmic Fairness}
\label{subsec:fairness}

Machine learning models trained on historical hospital data risk encoding
demographic biases present in the training distribution. Future work should
evaluate model performance disaggregated by age group, sex, ethnicity, and
imaging indication to identify potential disparities.

\subsection{Transparency and Auditability}
\label{subsec:transparency}

The PHI audit log (\code{phi\_audit.log}) provides a tamper-evident record
of every de-identification operation. The WandB training logs (when enabled)
capture all hyperparameters, loss curves, and validation metrics.
The discrepancy PDF reports generated by the clinical validation module
ensure that every model--radiologist disagreement is reviewed.
Together these mechanisms support the transparency requirements of the
EU AI Act Article~13 and the FDA's proposed framework for AI/ML-based
Software as a Medical Device~\citep{fda2021ai}.
